% 第3章3.4.3树基表示与节点势：五点示例（tree5.py 验算）。
% 供需 b=(4,0,-2,1,-3)；弧（容量，费用）：e1 1->2 (5,2)，e2 2->3 (4,1)，e3 4->3 (3,2)，e4 3->5 (5,1)，
% e5 1->4 (2,4) 属 U，e6 2->5 (3,3) 属 L。树流：e1=2,e2=2,e3=3（树弧恰在上界），e4=3。
% 以3为根：pi=(3,1,0,2,-1)；约化费用 e5: 4-3+2=3，e6: 3-1-1=1。
\begin{tikzpicture}[>=Stealth,line width=.8pt,font=\fontsize{8}{10}\selectfont,
  v/.style={circle,draw=FigureStroke,fill=FigureBlueFill,line width=.8pt,minimum size=6.5mm,inner sep=0pt,text=FigureInk},
  e/.style={->,draw=FigureStroke,line width=.8pt,shorten >=.5pt},
  tr/.style={->,draw=FigureBlue,line width=1.4pt,shorten >=.5pt},
  lo/.style={->,draw=FigureMuted,line width=.8pt,dashed,shorten >=.5pt},
  up/.style={->,draw=FigureMuted,line width=.9pt,densely dotted,shorten >=.5pt},
  lab/.style={text=FigureInk,inner sep=1.5pt},
  sup/.style={text=FigureBlue},
  note/.style={text=FigureMuted,font=\FigureNoteFont\fontsize{7.5}{9}\selectfont}]
  \foreach \dx/\p in {0/a,7.5/b}{
  \begin{scope}[xshift=\dx cm]
  \node[v] (\p1) at (0,0) {$1$};
  \node[v] (\p2) at (1.7,1) {$2$};
  \node[v] (\p3) at (3.4,0) {$3$};
  \node[v] (\p4) at (1.7,-1) {$4$};
  \node[v] (\p5) at (5.1,1) {$5$};
  \end{scope}}
  % (a) 树流
  \begin{scope}
  \draw[tr] (a1)--node[lab,above left=-1pt]{$2/5$} (a2);
  \draw[tr] (a2)--node[lab,above right=-1pt]{$2/4$} (a3);
  \draw[tr] (a4)--node[lab,below right=-1pt]{$3/3$} (a3);
  \draw[tr] (a3)--node[lab,below right=-1pt]{$3/5$} (a5);
  \draw[up] (a1)--node[lab,below left=-1pt]{$2/2$\,$U$} (a4);
  \draw[lo] (a2)--node[lab,above]{$0/3$\,$L$} (a5);
  \node[sup,left=1pt] at (a1.west) {$4$};
  \node[sup,above=1pt] at (a2.north) {$0$};
  \node[sup,right=1pt] at (a3.east) {$-2$};
  \node[sup,below=1pt] at (a4.south) {$1$};
  \node[sup,right=1pt] at (a5.east) {$-3$};
  \node at (2.55,-2.05) {(a) 固定非树弧后由叶求树流（$x/u$）};
  \end{scope}
  % (b) 节点势与约化费用
  \begin{scope}[xshift=7.5cm]
  \draw[tr] (b1)--node[lab,above left=-1pt]{$c=2$} (b2);
  \draw[tr] (b2)--node[lab,above right=-1pt]{$c=1$} (b3);
  \draw[tr] (b4)--node[lab,below right=-1pt]{$c=2$} (b3);
  \draw[tr] (b3)--node[lab,below right=-1pt]{$c=1$} (b5);
  \draw[up,draw=FigureRed] (b1)--node[lab,below left=-1pt,text=FigureRed]{$\bar c=3$} (b4);
  \draw[lo] (b2)--node[lab,above]{$\bar c=1$} (b5);
  \node[text=FigureGreen,left=1pt] at (b1.west) {$\pi=3$};
  \node[text=FigureGreen,above=1pt] at (b2.north) {$\pi=1$};
  \node[text=FigureGreen,right=1pt] at (b3.east) {$\pi=0$};
  \node[text=FigureGreen,below=1pt] at (b4.south) {$\pi=2$};
  \node[text=FigureGreen,right=1pt] at (b5.east) {$\pi=-1$};
  \node at (2.55,-2.05) {(b) 以$3$为根求势，再算非树弧约化费用};
  \end{scope}
  \node[note] at (6.3,-2.65) {粗蓝线：树弧$T$；虚线：取下界的$L$；点线：取上界的$U$；红色：违反最优性条件的非树弧};
\end{tikzpicture}
